% !PRINTED 575
\weberpage{575}
%<<CONT
Betrachtung in (3) $\mathfrak{m}$ jedes beliebige Ideal in $\Omega$ bedeuten kann, so kommen wir zu folgendem \gesperrt{Hauptsatz}:

\begin{quote}
% [sic] Satz 4 sagt im Original, T_m liege in Omega; Satz 3 beweist die Koeffizientenzugehoerigkeit nur zum Klassenkoerper. Diese unbegruendete Verengung des Grundkoerpers bleibt als Quellenbefund erhalten.
4. \gesperrt{Ist $\mathfrak{m}$ ein beliebiges Ideal des Körpers $\Omega$, so existiert eine Funktion $T_{\mathfrak m}$ in $\Omega$, deren Wurzeln die voneinander verschiedenen der Zahlen $\tau_\xi$ sind, wenn $\xi$ ein System inkongruenter, zu $\mathfrak{m}$ teilerfremder Zahlen durchläuft.}
\end{quote}

Um den Grad der Funktion $T_{\mathfrak m}$ zu bestimmen, bemerken wir, daß zwei Größen $\tau_\xi, \tau_{\xi'}$ nur dann einander gleich sind, wenn
\weq{8}{\xi \equiv \varepsilon\xi' \quad (\mathrm{mod}\ \mathfrak{m})}
ist. Es werden also so viele von den $\tau_\xi$ einander gleich, als es modulo $\mathfrak{m}$ inkongruente Einheiten gibt; das sind im allgemeinen zwei, für $\Delta = -4$ sind es vier und für $\Delta = -3$ sind es sechs. Diese Zahlen verringern sich aber wiederum, wenn verschiedene der Einheiten $\varepsilon$ nach dem Modul $\mathfrak{m}$ kongruent sind, wenn also $1 - \varepsilon$ durch $\mathfrak{m}$ teilbar ist. Das kann aber nur vorkommen, wenn $\mathfrak{m}$ ein Teiler von 2 oder von 3 ist. Also:

\begin{quote}
5. \gesperrt{Der Grad der Funktion $T_{\mathfrak m}$ ist gleich der Anzahl $\psi(\mathfrak{m})$} (Bd. II, \S\ 168) \gesperrt{der nach dem Modul $\mathfrak{m}$ inkongruenten, zu $\mathfrak{m}$ teilerfremden Zahlen in $\Omega$, geteilt durch die Anzahl der nach $\mathfrak{m}$ inkongruenten Einheiten in $\Omega$.}
\end{quote}

Die Wurzeln der Funktion $T_{\mathfrak m}$ bestimmen einen algebraischen Körper $\mathfrak{T}_{\mathfrak m}$ über dem Klassenkörper, dessen relativer Grad höchstens gleich dem Grade von $T_{\mathfrak m}$ ist, und beide Grade sind gleich, wenn $T_{\mathfrak m}$ irreduzibel ist.

\begin{quote}
6. \gesperrt{Diesen Körper $\mathfrak{T}_{\mathfrak m}$ wollen wir den Teilungskörper für den Modul $\mathfrak{m}$ nennen.}
\end{quote}

Diese Teilungskörper spielen für den Körper $\Omega$ eine ähnliche Rolle, wie die Kreisteilungskörper für den Körper der rationalen Zahlen, nur daß sich hier noch der Klassenkörper dazwischen schiebt, zu dem es im Körper der rationalen Zahlen, ebenso wie in jedem einklassigen Körper, kein Analogon gibt.

Nach dem Multiplikationstheorem (\S\ 151) kann
\weq{9}{\tau(\xi u) = \theta_\xi[\tau(u)]}
rational durch $\tau(u)$ ausgedrückt werden. Setzen wir $u = \xi'\dfrac{\alpha}{\mu}$, so folgt:
\weq{10}{\tau_{\xi\xi'} = \tau\left(\frac{\xi\xi'\alpha}{\mu}\right) = \theta_\xi(\tau_{\xi'}),}
%CONT>>
